\documentclass[10pt,a4paper]{amsart}
\usepackage{amsmath,amssymb,amsthm,graphicx}
\usepackage{verbatim}
\newenvironment{reference}{\comment}{\endcomment}
%\newenvironment{reference}{}{}
\newenvironment{slogan}{\comment}{\endcomment}
\newenvironment{history}{\comment}{\endcomment}

% For commutative diagrams we use Xy-pic
\usepackage[all]{xy}

% We use 2cell for 2-commutative diagrams.
\xyoption{2cell}
% Theorem environments.
%
\theoremstyle{plain}
\newtheorem{theorem}[subsection]{Theorem}
\newtheorem{proposition}[subsection]{Proposition}
\newtheorem{lemma}[subsection]{Lemma}

\theoremstyle{definition}
\newtheorem{definition}[subsection]{Definition}
\newtheorem{example}[subsection]{Example}
\newtheorem{exercise}[subsection]{Exercise}
\newtheorem{situation}[subsection]{Situation}

\theoremstyle{remark}
\newtheorem{remark}[subsection]{Remark}
\newtheorem{remarks}[subsection]{Remarks}

\numberwithin{equation}{subsection}

% Macros
%
\def\lim{\mathop{\mathrm{lim}}\nolimits}
\def\colim{\mathop{\mathrm{colim}}\nolimits}
\def\Spec{\mathop{\mathrm{Spec}}}
\def\Hom{\mathop{\mathrm{Hom}}\nolimits}
\def\Ext{\mathop{\mathrm{Ext}}\nolimits}
\def\SheafHom{\mathop{\mathcal{H}\!\mathit{om}}\nolimits}
\def\SheafExt{\mathop{\mathcal{E}\!\mathit{xt}}\nolimits}
\def\Sch{\mathit{Sch}}
\def\Mor{\mathop{\mathrm{Mor}}\nolimits}
\def\Ob{\mathop{\mathrm{Ob}}\nolimits}
\def\Sh{\mathop{\mathit{Sh}}\nolimits}
\def\NL{\mathop{N\!L}\nolimits}
\def\CH{\mathop{\mathrm{CH}}\nolimits}
\def\proetale{{pro\text{-}\acute{e}tale}}
\def\etale{{\acute{e}tale}}
\def\QCoh{\mathit{QCoh}}
\def\Ker{\mathop{\mathrm{Ker}}}
\def\Im{\mathop{\mathrm{Im}}}
\def\Coker{\mathop{\mathrm{Coker}}}
\def\Coim{\mathop{\mathrm{Coim}}}

% Boxtimes
%
\DeclareMathSymbol{\boxtimes}{\mathbin}{AMSa}{"02}

%
% Macros for moduli stacks/spaces
%
\def\QCohstack{\mathcal{QC}\!\mathit{oh}}
\def\Cohstack{\mathcal{C}\!\mathit{oh}}
\def\Spacesstack{\mathcal{S}\!\mathit{paces}}
\def\Quotfunctor{\mathrm{Quot}}
\def\Hilbfunctor{\mathrm{Hilb}}
\def\Curvesstack{\mathcal{C}\!\mathit{urves}}
\def\Polarizedstack{\mathcal{P}\!\mathit{olarized}}
\def\Complexesstack{\mathcal{C}\!\mathit{omplexes}}
% \Pic is the operator that assigns to X its picard group, usage \Pic(X)
% \Picardstack_{X/B} denotes the Picard stack of X over B
% \Picardfunctor_{X/B} denotes the Picard functor of X over B
\def\Pic{\mathop{\mathrm{Pic}}\nolimits}
\def\Picardstack{\mathcal{P}\!\mathit{ic}}
\def\Picardfunctor{\mathrm{Pic}}
\def\Deformationcategory{\mathcal{D}\!\mathit{ef}}
\usepackage{cleveref}
\crefname{originallemma}{original lemma}{original lemmas}
\Crefname{originallemma}{Original lemma}{Original lemmas}
\makeatletter
\newlabel{topology-lemma-dimension-function-unique}{{02IB}{source}{}{sourcepointer-1}{}}
\newlabel{topology-lemma-dimension-function-unique@cref}{{[originallemma][0][]02IB}{[1][0][]source}{}{}{}}
\makeatother
\crefname{originallemma}{original lemma}{original lemmas}
\Crefname{originallemma}{Original lemma}{Original lemmas}
\makeatletter
\newlabel{lemma-dimension-function-scheme}{{0AWF}{source}{}{sourcepointer-2}{}}
\newlabel{lemma-dimension-function-scheme@cref}{{[originallemma][0][]0AWF}{[1][0][]source}{}{}{}}
\makeatother
\crefname{originallemma}{original lemma}{original lemmas}
\Crefname{originallemma}{Original lemma}{Original lemmas}
\makeatletter
\newlabel{dualizing-lemma-depth-dualizing-module}{{0AWE}{source}{}{sourcepointer-3}{}}
\newlabel{dualizing-lemma-depth-dualizing-module@cref}{{[originallemma][0][]0AWE}{[1][0][]source}{}{}{}}
\makeatother
\crefname{originaldefinition}{original definition}{original definitions}
\Crefname{originaldefinition}{Original definition}{Original definitions}
\makeatletter
\newlabel{coherent-definition-depth}{{0341}{source}{}{sourcepointer-4}{}}
\newlabel{coherent-definition-depth@cref}{{[originaldefinition][0][]0341}{[1][0][]source}{}{}{}}
\makeatother
\title{The dualizing module of a connected Noetherian scheme is (S\_2) with support on the maximal dimensional components}
\author{The Stacks Project Authors}
\date{stacks-project:duality.tex:local-snapshot-20261012; GNU FDL 1.2 or later}
\begin{document}
\maketitle
\noindent\textit{Editorial note.} All delivered mathematical text is the author's own, copied from the pinned Stacks Project chapter file duality.tex; no statement or proof step was rewritten, reordered or removed.\par\medskip
\noindent\textit{Editorial note.} Preamble provenance: the pinned chapter file opens with an input of the shared preamble and carries no preamble text of its own, so the project's single Stacks-wide preamble staged at sources/stacks-duality/preamble.tex is excerpted verbatim under the recipe fixed by LANE\_BRIEF 22.1, as three spans, lines 8-13 (verbatim package and the reference/slogan/history environments), lines 14-18 (Xy-pic setup) and lines 154-217 (the native theorem declarations and mathematical macros). Nothing was added, removed or redefined; the wrapper class and the amsmath/amssymb/amsthm/graphicx base packages are supplied by the assembler.\par\medskip
\noindent\textit{Editorial note.} The counted statement and proof contain no cross reference, no citation and no included file, so the delivered body is self contained.\par\medskip
\noindent\textit{Editorial note.} Canonical location of the counted proposition: Stacks tag 0AWK (sources/stacks-duality/pinned-original-source.tex line 6101).\par\medskip
\noindent\textit{Allowed standard prerequisites (editorial).}\begin{itemize}
\item Uniqueness of a dimension function up to a constant on a connected scheme. Topology, Lemma labelled topology-lemma-dimension-function-unique, cited at proof lines 6128-6130. Allows choosing the dimension function of the dualizing complex.
\item The dimension function attached to a dualizing complex. Same chapter, Lemma labelled lemma-dimension-function-scheme, cited at proof lines 6130-6132. Supplies the function used.
\item The algebraic depth and support statement for the dualizing module. Dualizing Complexes, Lemma labelled dualizing-lemma-depth-dualizing-module, cited at proof lines 6134-6136, with Cohomology of Schemes, Definition labelled coherent-definition-depth, cited at proof lines 6136-6140. Gives both assertions.
\end{itemize}\medskip
\ifdefined\hypertarget\hypertarget{sourcepointer-1}{}\fi
\noindent\textit{Source reference.} Original lemma 02IB: sources/stacks-chow/tags.txt, source lines 3294--3294. Cross-chapter standard prerequisite located in the project tag list staged in the local snapshot; invoked at proof line 6113 of the counted proof. It is retained only as a navigation pointer; every load-bearing step is inside the delivered span.\par\medskip
\ifdefined\hypertarget\hypertarget{sourcepointer-2}{}\fi
\noindent\textit{Source reference.} Original lemma 0AWF: sources/stacks-duality/pinned-original-source.tex, source lines 567--567. Same-chapter pointer into the author environment lemma at source line 567; invoked at proof line 6116 of the counted proof. It is retained only as a navigation pointer; every load-bearing step is inside the delivered span.\par\medskip
\ifdefined\hypertarget\hypertarget{sourcepointer-3}{}\fi
\noindent\textit{Source reference.} Original lemma 0AWE: sources/stacks-chow/tags.txt, source lines 13717--13717. Cross-chapter standard prerequisite located in the project tag list staged in the local snapshot; invoked at proof line 6119 of the counted proof. It is retained only as a navigation pointer; every load-bearing step is inside the delivered span.\par\medskip
\ifdefined\hypertarget\hypertarget{sourcepointer-4}{}\fi
\noindent\textit{Source reference.} Original definition 0341: sources/stacks-chow/tags.txt, source lines 4076--4076. Cross-chapter standard prerequisite located in the project tag list staged in the local snapshot; invoked at proof line 6121 of the counted proof. It is retained only as a navigation pointer; every load-bearing step is inside the delivered span.\par\medskip
\begin{lemma}
\label{lemma-dualizing-module}
Let $X$ be a connected Noetherian scheme and let $\omega_X$ be a dualizing
module on $X$. The support of $\omega_X$ is the union of the irreducible
components of maximal dimension with respect to any dimension function
and $\omega_X$ is a coherent $\mathcal{O}_X$-module having property $(S_2)$.
\end{lemma}

\begin{proof}
By our conventions discussed above there exists a dualizing complex
$\omega_X^\bullet$ such that $\omega_X$ is the leftmost nonvanishing
cohomology sheaf. Since $X$ is connected, any two dimension functions
differ by a constant
(Topology, Lemma \ref{topology-lemma-dimension-function-unique}).
Hence we may use the
dimension function associated to $\omega_X^\bullet$
(Lemma \ref{lemma-dimension-function-scheme}).
With these remarks in place, the lemma now
follows from Dualizing Complexes, Lemma
\ref{dualizing-lemma-depth-dualizing-module}
and the definitions (in particular
Cohomology of Schemes, Definition \ref{coherent-definition-depth}).
\end{proof}

\end{document}
